
\medskip
\noindent\textbf{Proposition 2.3.1.}
\phantomsection
\label{sga3:VIIB:2.3.1}
{\itshape Let \(K\) and \(G\) be two formal \(k\)-groups, with
affine algebras \(B\) and \(A\), respectively. Suppose that \(K\) is
topologically flat over \(k\). Then there is a canonical bijection
from
\[
  \operatorname{Hom}_{\operatorname{Grf}/k}(K,G)
\]
onto the set of homomorphisms of unital \(\mathbf O_k\)-algebras
\[
  h:\mathbf H(K)\longrightarrow\mathbf H(G)
\]
for which the diagram
\begin{center}
  \phantomsection
  \label{sga3:VIIB:2.3.1:compatibility-diagram}
  \begin{tikzcd}[row sep=1.25em,column sep=4em]
\mathbf H(K)\otimes_k\mathbf H(K)
  \arrow[r,"h\otimes h"]
  & \mathbf H(G)\otimes_k\mathbf H(G)
      \arrow[dr,"\varphi_G"] & {} \\
{} & {} & \mathbf H(G\times G) \\
\mathbf H(K) \arrow[uu,"\Delta_{\mathbf H(K)}"]
  \arrow[r,"h"]
  & \mathbf H(G) \arrow[ur,"\delta_G"']
  & {}
\end{tikzcd}

\end{center}
is commutative.}

Since \(K\) is topologically flat, \(\mathbf H(K)\) is endowed with
a bialgebra structure \textup{(cf.\ \SGARefLink{sga3:VIIB:2.2}{2.2})}, and
\(\Delta_{\mathbf H(K)}\) is its diagonal morphism. In other words,
with the notation of \SGARefLink{sga3:VIIB:2.3}{2.3},
\[
  \Delta_{\mathbf H(K)}=\phi_K^{-1}\circ\delta_K.
\]
When \(G\) too is topologically flat over \(k\), the proposition
follows from the equivalence of categories established in \SGARefLink{sga3:VIIB:2.2.1}{2.2.1}.

In the general case, we may suppose that \(k\) is Artinian and argue
with the algebras
\[
  \mathbf H(K)=B^\dagger,
  \qquad
  \mathbf H(G)=A^\dagger.
\]
Let \(\operatorname{Hom}_c(A,B)\) denote the set of continuous
\(k\)-linear maps from \(A\) to \(B\), and let
\(\operatorname{Hom}_k(B^\dagger,A^\dagger)\) denote the set of
\(k\)-linear maps from \(B^\dagger\) to \(A^\dagger\).
\footnote{\label{sga3:VIIB:nde:87}The editors supplied the details of
the remainder of the proof.}

By \SGARefLink{sga3:VIIB:0.3.6.A}{0.3.6.A}, if \(M\) and \(P\) are pseudocompact \(k\)-modules and
\(P\) is projective, the canonical map
\[
  \operatorname{Hom}_c(M,P)
  \longrightarrow
  \operatorname{Hom}_k(P^\dagger,M^\dagger),
  \qquad
  f\longmapsto{}^t f,
\]
where \({}^t f\) denotes the transpose of \(f\), is bijective. We
shall apply this either with
\[
  M=A\widehat\otimes A,\quad P=B,
\]
or with
\[
  M=A,\quad P=B\widehat\otimes B.
\]

Let \(f\in\operatorname{Hom}_c(A,B)\). Consider the following
diagrams. The squares marked \textup{\SGARefLink{sga3:VIIB:0}{(0)}} are commutative, and the
two unnamed vertical arrows in the second diagram are
\({}^t(f\widehat\otimes f)\).
\begin{center}
  \phantomsection
  \label{sga3:VIIB:2.3.1:affine-diagram}
  \begin{tikzcd}[row sep=3.1em,column sep=3.5em]
A\widehat\otimes A
  \arrow[r,"m_A"] \arrow[d,"{f\widehat\otimes f}"']
  \arrow[dr,phantom,"\textup{(1)}" description]
  & A \arrow[r,"\Delta_A"] \arrow[d,"f"]
      \arrow[dr,phantom,"\textup{(2)}" description]
  & A\widehat\otimes A \arrow[d,"{f\widehat\otimes f}"] \\
B\widehat\otimes B \arrow[r,"m_B"]
  & B \arrow[r,"\Delta_B"]
  & B\widehat\otimes B
\end{tikzcd}


  \medskip
  \phantomsection
  \label{sga3:VIIB:2.3.1:dual-diagram}
  \begin{tikzcd}[row sep=3.4em,column sep=2.15em]
A^\dagger\otimes A^\dagger
  \arrow[r,"\varphi_G"]
  \arrow[dr,phantom,"\textup{(0)}" description]
  & (A\widehat\otimes A)^\dagger
  & A^\dagger
      \arrow[l,"{\delta_G=m_A^t}"']
      \arrow[dr,phantom,"\textup{(2')}" description]
  & (A\widehat\otimes A)^\dagger
      \arrow[l,"\Delta_A^t"']
      \arrow[dr,phantom,"\textup{(0)}" description]
  & A^\dagger\otimes A^\dagger \arrow[l,"\varphi_G"'] \\
B^\dagger\otimes B^\dagger
  \arrow[r,equals] \arrow[u,"{{}^tf\otimes{}^tf}"]
  & (B\widehat\otimes B)^\dagger
      \arrow[u]
  & B^\dagger
      \arrow[l,"{\delta_K=m_B^t}"]
      \arrow[u,"{{}^tf}"']
      \arrow[ul,phantom,"\textup{(1')}" description]
  & (B\widehat\otimes B)^\dagger
      \arrow[l,"\Delta_B^t"]
      \arrow[u]
  & B^\dagger\otimes B^\dagger
      \arrow[l,equals] \arrow[u,"{{}^tf\otimes{}^tf}"']
\end{tikzcd}

\end{center}

If \(f:A\to B\) corresponds to a morphism of formal groups
\(K\to G\), then the squares \textup{(1)} and \textup{(2)} are
commutative and
\[
  \varepsilon_B\circ f=\varepsilon_A.
\]
Consequently, the squares \textup{(1\(^{\prime}\))} and
\textup{(2\(^{\prime}\))} are commutative, and \({}^t f\) sends the
unit of \(B^\dagger=\mathbf H(K)\) to that of
\(A^\dagger=\mathbf H(G)\). Thus \({}^t f\) is a homomorphism of
unital \(k\)-algebras
\[
  \mathbf H(K)\longrightarrow\mathbf H(G)
\]
for which the diagram in the proposition is commutative.

Conversely, suppose that \({}^t f\) satisfies these conditions.
Then
\[
  \varepsilon_B\circ f=\varepsilon_A
\]
and the squares \textup{(1\(^{\prime}\))} and
\textup{(2\(^{\prime}\))} are commutative. For
\[
  M=A\widehat\otimes A,\quad P=B
  \qquad
  \text{(respectively \(M=A,\ P=B\widehat\otimes B\))},
\]
the map \(g\mapsto{}^t g\) is injective. Hence the squares
\textup{(1)} and \textup{(2)} are commutative, so \(f\) is
compatible with the multiplications and diagonal morphisms of \(A\)
and \(B\).

It remains to prove that \(f(1_A)=1_B\). What precedes gives
\[
  \varepsilon_B f(1)=1,
  \qquad
  \Delta_B f(1)=f(1)\widehat\otimes f(1),
  \qquad
  f(1)\cdot f(1)=f(1).
\]
By \SGARefLink{sga3:VIIA:2.1}{2.1}\textup{(iii)}, the first two conditions imply that \(f(1)\)
has inverse \(c_Bf(1)\) in \(B\). The equality
\(f(1)\cdot f(1)=f(1)\) therefore implies \(f(1)=1\). Thus
\[
  f:A\longrightarrow B
\]
is a morphism in \(\operatorname{Alp}/k\), compatible with the
comultiplications of \(A\) and \(B\).
